\input{../../../templates/es/preambulo.tex}

% Los listados se toman de src/ con \lstinputlisting; la diapositiva es el archivo.
\lstset{
    basicstyle=\ttfamily\fontsize{8pt}{9.6pt}\selectfont,
    breaklines=true,
    breakatwhitespace=true,
    columns=fullflexible,
    keepspaces=true,
    xleftmargin=0.3em,
    framesep=2pt,
    aboveskip=0pt,
    belowskip=0pt
}

\newcommand{\marcoresultado}[3]{%
    \begin{frame}{#1}
        \begin{center}
            \includegraphics[height=0.62\textheight,width=\textwidth,keepaspectratio]{#2}
        \end{center}
        \vspace{0.25em}
        {\small #3\par}
    \end{frame}%
}

\newcommand{\marcoafirmacion}[2]{%
    \begin{frame}{#1}
        {\small #2\par}
    \end{frame}%
}

\date{}

\begin{document}

\tituloleccion{13}{Mejorando el Rendimiento}

\section{Esta lección}

\begin{frame}{Esta lección}
El censo de evaluaciones de la Lección 02 era un número sin método. Esta lección es el método: dos contadores, una plantilla, siete pasos, la misma respuesta.

\vspace{0.6em}
\begin{itemize}
    \item Evaluar una vez y guardar: misma semilla, mismo \(-87\), 54\% menos llamadas.
    \item Una caché vale lo que la búsqueda repite --- medido, no supuesto.
    \item Solo un objetivo más grueso achica de verdad el problema, y hay que juzgarlo en el real.
\end{itemize}
\end{frame}

% ================= turnos_01_el_costo_de_una_ejecucion =================
\begin{frame}{Un modelo de trabajo}
\[
T\approx N_{\mathrm{llamadas}}W_{\mathrm{llamada}}+T_{\mathrm{serial}}
+T_{\mathrm{procesos}}.
\]
Una caché exige un objetivo determinista indexado por el genoma inmutable
completo. Un reinicio exacto también necesita población, estado de
aptitudes/caché, contadores de generación y estado
pseudoaleatorio. Los workers tienen memoria separada con spawn de Windows y con
fork de copia al escribir en Unix; sus cachés y contadores no se comparten solos.
\end{frame}

\section{El costo real de una ejecución}

\begin{frame}{El costo real de una ejecución}
Esta es la última lección del curso, y es la única sobre el costo de
ejecutar un algoritmo genético en lugar de sobre el algoritmo en sí. Todo lo
que sigue es un intento de hacer que esta ejecución sea más barata. Así que lo primero que hay que construir
no es una optimización: es el medidor.
\end{frame}

\marcoresultado{El costo real de una ejecución}{../figuras/turnos_01_el_costo_de_una_ejecucion.png}{%
    La ejecución construye \textbf{428} individuos y llama a la aptitud \textbf{930} veces — \textbf{2.17 evaluaciones por individuo}, dividido en SELECCIÓN 300 / REPRODUCCIÓN 0 / ESTADÍSTICAS 630. \textbf{195,300} unidades de trabajo, mejor aptitud \textbf{−87}}

% ================= turnos_02_evaluar_una_vez =================
\section{Evaluar la función de aptitud una vez}

\begin{frame}{Evaluar la función de aptitud una vez}
Esta es la sección 13.1 del libro, y es el mejor trato de todo el curso.
Los genes de un individuo no cambian después de que se construye, por lo que su aptitud tampoco puede
cambiar. A pesar de esto, el paso 1 lo recalculaba cada vez que se necesitaba un valor
- una vez por individuo dentro del ordenamiento, dos veces más por individuo en las
estadísticas - y pagaba una evaluación completa de 210 celdas cada vez.
\end{frame}

\begin{frame}[fragile]{(1) Individuo}
\lstinputlisting[basicstyle=\ttfamily\fontsize{8pt}{9.6pt}\selectfont,firstline=102,lastline=115]{../src/turnos_02_evaluar_una_vez.py}
\end{frame}

\begin{frame}[fragile]{(2) aptitud\_de()}
\lstinputlisting[basicstyle=\ttfamily\fontsize{8pt}{9.6pt}\selectfont,firstline=118,lastline=122]{../src/turnos_02_evaluar_una_vez.py}
\end{frame}

\begin{frame}[fragile]{(3) el\_reporte}
\lstinputlisting[basicstyle=\ttfamily\fontsize{6pt}{7.2pt}\selectfont,firstline=220,lastline=288]{../src/turnos_02_evaluar_una_vez.py}
\end{frame}

\marcoresultado{Evaluar la función de aptitud una vez}{../figuras/turnos_02_evaluar_una_vez.png}{%
    \textbf{428} llamadas, \textbf{89,880} unidades de trabajo: \textbf{54.0\%} menos, con SELECCIÓN y ESTADÍSTICAS en \textbf{0}. La mejor aptitud sigue siendo \textbf{−87} — el ahorro no costó nada}

% ================= turnos_03_la_cache =================
\section{Una caché para genomas que regresan}

\begin{frame}{Una caché para genomas que regresan}
El paso 2 puso la factura en una evaluación por individuo construido. Ese es un piso
solo para genomas *nuevos*. Un algoritmo genético no produce genomas nuevos todo
el tiempo: las élites se copian hacia adelante, la selección elige a los mismos padres dos veces,
la cruza entre dos copias de una asignación devuelve esa asignación sin cambios, y
una mutación puede deshacer una anterior. Cada uno de esos es un genoma que el programa
ya ha cotizado.

(h=H/(H+M)). Si domina el trabajo del objetivo, la aceleración ideal por
trabajo es (1/(1-h)); la observada también paga búsqueda y procesos.
\end{frame}

\begin{frame}[fragile]{(1) CACHE}
\lstinputlisting[basicstyle=\ttfamily\fontsize{8pt}{9.6pt}\selectfont,firstline=97,lastline=103]{../src/turnos_03_la_cache.py}
\end{frame}

\begin{frame}[fragile]{(2) aptitud\_en\_cache()}
\lstinputlisting[basicstyle=\ttfamily\fontsize{8pt}{9.6pt}\selectfont,firstline=112,lastline=126]{../src/turnos_03_la_cache.py}
\end{frame}

\begin{frame}[fragile]{(3) Individuo}
\lstinputlisting[basicstyle=\ttfamily\fontsize{8pt}{9.6pt}\selectfont,firstline=139,lastline=139]{../src/turnos_03_la_cache.py}
\end{frame}

\begin{frame}[fragile]{(4) el\_reporte}
\lstinputlisting[basicstyle=\ttfamily\fontsize{6pt}{7.2pt}\selectfont,firstline=244,lastline=310]{../src/turnos_03_la_cache.py}
\end{frame}

\begin{frame}[fragile]{(5) escaneo\_tasa\_aciertos()}
\lstinputlisting[basicstyle=\ttfamily\fontsize{6pt}{7.2pt}\selectfont,firstline=313,lastline=351]{../src/turnos_03_la_cache.py}
\end{frame}

\marcoafirmacion{Una caché para genomas que regresan}{%
    \textbf{19} aciertos en 428 construcciones, \textbf{409} genomas distintos, 4.4\% menos trabajo, respuesta sin cambios en \textbf{−87}. El escaneo muestra que la tasa de aciertos es propiedad de la *búsqueda*: \textbf{14.4\%} con mutación p=0.10, \textbf{6.1\%} a 0.25, \textbf{4.4\%} a 0.50, \textbf{3.2\%} a 0.75}

% ================= turnos_04_instantanea =================
\section{Instantáneas, y lo que olvida una instantánea}

\begin{frame}{Instantáneas, y lo que olvida una instantánea}
Las ejecuciones largas se interrumpen: una laptop se suspende, un trabajo en cola alcanza su tiempo límite, un
estudio de afinación quiere ramificarse de una ejecución que ya costó una hora. La respuesta del libro
(sección 13.4) es escribir los genes de la población en un archivo y leerlos
de nuevo más tarde.
\end{frame}

\begin{frame}[fragile]{(1) INSTANTANEAS}
\lstinputlisting[basicstyle=\ttfamily\fontsize{8pt}{9.6pt}\selectfont,firstline=60,lastline=64]{../src/turnos_04_instantanea.py}
\end{frame}

\begin{frame}[fragile]{(2) volcar\_poblacion()}
\lstinputlisting[basicstyle=\ttfamily\fontsize{8pt}{9.6pt}\selectfont,firstline=152,lastline=158]{../src/turnos_04_instantanea.py}
\end{frame}

\begin{frame}[fragile]{(3) restaurar\_poblacion()}
\lstinputlisting[basicstyle=\ttfamily\fontsize{8pt}{9.6pt}\selectfont,firstline=161,lastline=166]{../src/turnos_04_instantanea.py}
\end{frame}

\begin{frame}[fragile]{(4) volcar\_cache()}
\lstinputlisting[basicstyle=\ttfamily\fontsize{8pt}{9.6pt}\selectfont,firstline=169,lastline=175]{../src/turnos_04_instantanea.py}
\end{frame}

\begin{frame}[fragile]{(5) restaurar\_cache()}
\lstinputlisting[basicstyle=\ttfamily\fontsize{8pt}{9.6pt}\selectfont,firstline=178,lastline=184]{../src/turnos_04_instantanea.py}
\end{frame}

\begin{frame}[fragile]{(6) el\_reporte}
\lstinputlisting[basicstyle=\ttfamily\fontsize{6pt}{7.2pt}\selectfont,firstline=283,lastline=382]{../src/turnos_04_instantanea.py}
\end{frame}

\marcoafirmacion{Instantáneas, y lo que olvida una instantánea}{%
    30 asignaciones se restauran por \textbf{30 evaluaciones / 6,300 unidades de trabajo} a partir de genes solos, y por \textbf{0 / 0} cuando la caché se restaura primero. El archivo de caché es \textbf{14.0×} el tamaño del archivo de genes — un mal intercambio a 210 unidades por llamada, uno excelente a un minuto por llamada}

% ================= turnos_05_evaluacion_paralela =================
\section{Evaluando una generación en un pool de procesos}

\begin{frame}{Evaluando una generación en un pool de procesos}
Esta es la sección 13.5 del libro, y la primera técnica en esta lección que
no elimina ningún trabajo en absoluto. La caché elimina evaluaciones. Un pool no elimina
ninguna: ejecuta las mismas evaluaciones en otro lugar, al mismo tiempo.
\end{frame}

\begin{frame}[fragile]{(1) Individuo}
\lstinputlisting[basicstyle=\ttfamily\fontsize{8pt}{9.6pt}\selectfont,firstline=155,lastline=169]{../src/turnos_05_evaluacion_paralela.py}
\end{frame}

\begin{frame}[fragile]{(2) evaluar\_lote()}
\lstinputlisting[basicstyle=\ttfamily\fontsize{8pt}{9.6pt}\selectfont,firstline=172,lastline=182]{../src/turnos_05_evaluacion_paralela.py}
\end{frame}

\begin{frame}[fragile]{(3) materializar()}
\lstinputlisting[basicstyle=\ttfamily\fontsize{8pt}{9.6pt}\selectfont,firstline=185,lastline=201]{../src/turnos_05_evaluacion_paralela.py}
\end{frame}

\begin{frame}[fragile]{(4) operacion\_cruza()}
\lstinputlisting[basicstyle=\ttfamily\fontsize{8pt}{9.6pt}\selectfont,firstline=274,lastline=285]{../src/turnos_05_evaluacion_paralela.py}
\end{frame}

\begin{frame}[fragile]{(5) operacion\_mutacion()}
\lstinputlisting[basicstyle=\ttfamily\fontsize{8pt}{9.6pt}\selectfont,firstline=288,lastline=299]{../src/turnos_05_evaluacion_paralela.py}
\end{frame}

\begin{frame}[fragile]{(6) ejecutar()}
\lstinputlisting[basicstyle=\ttfamily\fontsize{7pt}{8.5pt}\selectfont,firstline=311,lastline=338]{../src/turnos_05_evaluacion_paralela.py}
\end{frame}

\begin{frame}[fragile]{(7) el\_reporte}
\lstinputlisting[basicstyle=\ttfamily\fontsize{6pt}{7.2pt}\selectfont,firstline=344,lastline=481]{../src/turnos_05_evaluacion_paralela.py}
\end{frame}

\marcoafirmacion{Evaluando una generación en un pool de procesos}{%
    El solo agrupamiento reduce la ejecución de 409 a \textbf{291} evaluaciones (\textbf{−28.9\%}), porque un hijo cruzado-luego-mutado solía evaluarse dos veces. La ejecución agrupada da la misma respuesta y el mismo mejor-hasta-ahora en las 10 generaciones — mientras que los contadores del padre leen \textbf{0 evaluaciones, 0 unidades de trabajo}}

% ================= turnos_06_cache_entre_procesos =================
\section{Llevando la caché y el medidor de vuelta a través de la frontera}

\begin{frame}{Llevando la caché y el medidor de vuelta a través de la frontera}
El paso 5 terminó con una ejecución agrupada cuyos contadores leían cero evaluaciones y cero
unidades de trabajo, para una ejecución que claramente había hecho ambas cosas.
Un worker tiene memoria separada: Windows usa spawn y Unix inicia con fork y
copia al escribir. En ambos casos, su caché y sus contadores no se comparten
automáticamente con el proceso padre.
\end{frame}

\begin{frame}[fragile]{(1) evaluar\_con\_costo()}
\lstinputlisting[basicstyle=\ttfamily\fontsize{8pt}{9.6pt}\selectfont,firstline=137,lastline=148]{../src/turnos_06_cache_entre_procesos.py}
\end{frame}

\begin{frame}[fragile]{(2) evaluar\_lote()}
\lstinputlisting[basicstyle=\ttfamily\fontsize{6pt}{7.2pt}\selectfont,firstline=151,lastline=184]{../src/turnos_06_cache_entre_procesos.py}
\end{frame}

\begin{frame}[fragile]{(3) el\_reporte}
\lstinputlisting[basicstyle=\ttfamily\fontsize{6pt}{7.2pt}\selectfont,firstline=338,lastline=417]{../src/turnos_06_cache_entre_procesos.py}
\end{frame}

\marcoafirmacion{Llevando la caché y el medidor de vuelta a través de la frontera}{%
    La ejecución agrupada ahora reporta \textbf{291 / 8 aciertos / 61,110}, igualando a la ejecución secuencial en los tres, respuesta sigue en \textbf{−87}. Las cachés privadas por trabajador del paso 5 descartaban todos los aciertos}

% ================= turnos_07_aptitud_gruesa =================
\section{Una aptitud gruesa, y lo que cuesta ser barato}

\begin{frame}{Una aptitud gruesa, y lo que cuesta ser barato}
Cada optimización hasta ahora ha sido gratuita. La sección 13.1 dejó de recalcular un
número conocido; la caché dejó de recalcular un genoma conocido; el agrupamiento dejó de
evaluar hijos que estaban a punto de ser descartados; el pool movió el trabajo
sin cambiarlo. Ninguna de ellas alteró un solo valor de aptitud, y cada
una se verificó contra la respuesta del paso anterior para demostrarlo.

Para el sustituto \(\tilde f\), \(e(x)=\tilde f(x)-f(x)\). La selección puede
usar \(\tilde f\), pero los candidatos finales deben calificarse con el \(f\) real.
\end{frame}

\begin{frame}[fragile]{(1) aptitud\_gruesa()}
\lstinputlisting[basicstyle=\ttfamily\fontsize{8pt}{9.6pt}\selectfont,firstline=138,lastline=149]{../src/turnos_07_aptitud_gruesa.py}
\end{frame}

\begin{frame}[fragile]{(2) evaluar\_con\_costo()}
\lstinputlisting[basicstyle=\ttfamily\fontsize{8pt}{9.6pt}\selectfont,firstline=152,lastline=164]{../src/turnos_07_aptitud_gruesa.py}
\end{frame}

\begin{frame}[fragile]{(3) el\_reporte}
\lstinputlisting[basicstyle=\ttfamily\fontsize{6pt}{7.2pt}\selectfont,firstline=352,lastline=448]{../src/turnos_07_aptitud_gruesa.py}
\end{frame}

\marcoafirmacion{Una aptitud gruesa, y lo que cuesta ser barato}{%
    \textbf{105} unidades de trabajo por llamada en lugar de 210, \textbf{51.0\%} menos trabajo para la ejecución. La búsqueda gruesa llega a \textbf{0 en su propio objetivo} — y su asignación puntúa \textbf{−160} en el real frente a \textbf{−87}, con \textbf{32} violaciones de descanso en lugar de 17}

\section{Siguiente}

\begin{frame}{Siguiente}
Los pasos 1--6 reportan todos la mejor aptitud \(-87\). El contenido es la factura que baja debajo. La aptitud gruesa da 0 en sí misma y \(-160\) en el objetivo real.

\vspace{0.8em}
\textbf{El reto }

Aplicar la misma disciplina: declarar el modelo matemático, mantener visible la factibilidad, comparar con una sola moneda de evaluación y separar desempeño observado de prueba. La especificación vigente vive en el repositorio del semestre.
\end{frame}
\end{document}
